\section{Setup and strategy of the proofs}
\label{sec:general-setup}

In this section, we recall the basic terminology and describe the strategies used in the proof of the main theorem.
We first review basic notions and properties of numerical semigroups and numerical semigroup rings, and then describe the strategies used to prove the Cohen--Macaulayness and normality of the Rees algebras appearing in the main theorem.

In Subsections~\ref{subsec:cm-strategy} and \ref{subsec:normality}, whenever
an infinite ground field is needed, we may replace $k$ by a purely
transcendental extension and assume that it is infinite; the assertions
concerning Cohen--Macaulayness and integral closedness descend by faithful
flatness.

\subsection{Basic notions of numerical semigroups and numerical semigroup rings}\label{subsec:basicnotion}

A \defword{numerical semigroup} $H$ is a submonoid of $\mathbb{N}$ with finite complement $\mathbb{N} \setminus H$.
For every numerical semigroup $H$, there exists a unique minimal system of generators $\{a_1, a_2, \ldots , a_n\}$ of $H$; that is, $H=\{\lambda_1 a_1 + \lambda_2 a_2 + \cdots + \lambda_n a_n \mid \lambda_1, \lambda_2, \ldots , \lambda_n\in \mathbb{N}\}$, and no proper subset of $\{a_1,a_2,\ldots,a_n\}$ generates $H$. Note that the greatest common divisor of $a_1, a_2, \ldots , a_n$ must be one. 
We write $\left<a_1, a_2, \ldots , a_n\right>$ for the numerical semigroup generated by $a_1, a_2 ,\ldots , a_n$.
The number $n$ is called the \defword{embedding dimension} of $H$, denoted by $\embdim(H)$, and the least positive integer in $H$ is called the \defword{multiplicity} of $H$, denoted by $\mult(H)$.
We have $\embdim(H)\leq\mult(H)$.
For any $h \in H \setminus \{0\}$, we put
$$
\Ap(H,h) = \{h' \in H \mid h' - h \notin H\}
$$
and call it the \defword{Ap\'ery set} of $H$ with respect to $h$. Obviously, if we put
$$
h'_i = \min \{h' \in H \mid h' \equiv i \pmod{h}\}
$$
for each $0 \le i \le h-1$, then $\Ap(H,h) = \{h'_0, h'_1, \ldots , h'_{h-1}\}$.
The largest integer not belonging to $H$ is called the \defword{Frobenius number} of $H$, denoted by $\Fr(H)$.
Hence, we have $a \in H$ for every integer $a \ge \Fr(H) + 1$, and $\Fr(H) = \max \Ap(H,h) - h$ for every $h \in H \setminus \{0\}$. 

Throughout this paper, let $k$ denote an arbitrary field. The semigroup ring of $H = \left<a_1, a_2, \ldots , a_n\right>$ over $k$ is the $k$-subalgebra 
$$
k[H] = k[t^h \mid h \in H] = k[t^{a_1}, t^{a_2}, \ldots , t^{a_n}]
$$
of the polynomial ring $k[t]$. The ring $k[H]$ is a one-dimensional Cohen--Macaulay domain. We regard $k[t]$ as a graded ring by $\deg t =1$, and hence $k[H]$ is a graded subring of $k[t]$.
The embedding dimension of $H$ coincides with the embedding dimension of the ring $k[H]$, and the multiplicity of $H$ is equal to the multiplicity of $k[H]$ with respect to the irrelevant maximal ideal $\fkm = (t^{a_1}, t^{a_2}, \ldots , t^{a_n})$.
Note that, for every $h \in H\setminus \{0\}$, the residue classes of the monomials $t^{h'}$, where $h'\in\Ap(H,h)$, form a basis of the $k$-vector space $k[H]/t^hk[H]$. 
Hence $k[H]$ is Gorenstein if and only if $\Ap(H,h)$ is symmetric in the following sense for every (equivalently, for some) $h\in H \setminus \{0\}$: if we write $\Ap(H,h) = \{c_1 < c_2 < \cdots < c_h\}$, then $c_i + c_{h+1-i} = c_h$ for each $1 \le i \le h$. When this is the case, we call $H$ a \defword{symmetric} numerical semigroup.
Suppose in addition that $H\ne\mathbb N$. If $\embdim(H)=\mult(H)$, equivalently, if the ring $k[H]$ has maximal embedding dimension, then $H$ is symmetric if and only if $\embdim(H) = 2$.

Let $H=\left<a_1, a_2, \ldots ,a_n\right>$, and let $S=k[x_1, x_2, \ldots , x_n]$ be the polynomial ring graded by $\deg x_i =a_i$.
Consider the graded ring homomorphism $\varphi_H: S \to k[H]$ defined by $\varphi_H(x_i) = t^{a_i}$, and put $\fkp_H = \Ker \varphi_H$. We call $\fkp_H$ the \defword{defining ideal} of $k[H]$.

Since only the case $n=3$ will be discussed in the remainder of this paper, the basic results for this case are summarized below. The following results are due to Herzog \cite{Herzog1970}.


\begin{theorem}[\cite{Herzog1970}]\label{thm:herzog}
	Let $H=\left<a_1, a_2, a_3\right>$ be a numerical semigroup with $\embdim(H) = 3$. Then the following assertions hold.
	\begin{enumerate}
		\item $H$ is symmetric if and only if, after a suitable reordering of $a_1, a_2, a_3$, the greatest common divisor $d$ of $a_1$ and $a_2$ is at least 2 and $a_3 \in \left<\frac{a_1}{d}, \frac{a_2}{d}\right>$.
		\item For each $1\le i\le3$, put $q_i = \min\{q >0\mid qa_i \in \left<\{a_1, a_2,a_3\} \setminus \{a_i\}\right>\}$ and choose any $p_{ij}\in\mathbb N$ $(j\ne i)$ such that $q_i a_i = \sum_{j\ne i}p_{ij}a_j$. Then $$\fkp_H = \left(x_i^{q_i} - \prod_{j\ne i}x_j^{p_{ij}} ~\middle|~ i = 1,2,3\right).$$
		If $H$ is non-symmetric, these three binomials form a minimal system of
		generators of $\fkp_H$.
		\item If $H$ is symmetric and $a_1,a_2,a_3$ have been ordered as in (1), then $$\fkp_H = (x_1^{\frac{a_2}{d}} - x_2^{\frac{a_1}{d}}, x_3^d - x_1^{p_1}x_2^{p_2}),$$ where $a_3 = p_1 \frac{a_1}{d} + p_2 \frac{a_2}{d}$ $(p_1, p_2 \in \mathbb{N})$. These two binomials form a minimal system of generators of $\fkp_H$, and this agrees with the presentation in (2).
	\end{enumerate}
\end{theorem}

\begin{example}
	\begin{enumerate}
		\item Let $H=\left<3,5,7\right>$. It is not symmetric because $3,5,7$ are pairwise coprime. In this case, $q_1 = 4$, $q_2 = q_3 = 2$, $p_{12} = p_{13} = p_{21} = p_{23} = p_{32} = 1$, and $p_{31} = 3$. Hence $\fkp_H= (x_1^4 - x_2x_3, x_2^2-x_1x_3,x_3^2-x_1^3x_2)$.
		\item Let $H=\left<4,6,11\right>$. Since $11 = 4 \cdot \frac{4}{2} + \frac{6}{2} \in \left<\frac{4}{2}, \frac{6}{2}\right>$, $H$ is symmetric, and $\fkp_H = (x_1^3-x_2^2, x_3^2-x_1^4x_2)$.
	\end{enumerate}
\end{example}

\begin{corollary}\label{cor:emb4symm}
	Let $H=\left<4,a_2, a_3\right>$ be a numerical semigroup with $\embdim(H) = 3$ and $\mult(H) = 4$.
	Then $H$ is symmetric if and only if one of $a_2$ and $a_3$ is even. When this is the case, the other generator must be odd.
	Moreover, there exist an odd integer $p\geq3$ and an integer $q\geq1$ such
	that $H=\left<4,2p,p+2q\right>$.
\end{corollary}


\begin{proof}
	Suppose that $H$ is symmetric and that both $a_2$ and $a_3$ are odd. 
	Then, by Theorem~\ref{thm:herzog} (1), $d = \gcd(a_2, a_3) \ge 2$ and $4 \in \left<\frac{a_2}{d}, \frac{a_3}{d}\right>$. For each $i=2,3$, $\frac{a_i}{d}$ cannot equal $1$, since otherwise $\embdim(H) < 3$. 
	Since both $\frac{a_2}{d}$ and $\frac{a_3}{d}$ are greater than one, an expression of $4$ as a non-negative integral combination of them forces one of them to be either $2$ or $4$. 
	Hence $\frac{a_i}{d}$ must be even for some $i=2,3$, which implies that $a_i$ is even; this is a contradiction.

	Conversely, suppose that one of $a_2$ and $a_3$ is even.  Denote the even
	generator by $2p$ and the other generator by $a$.  Then $a$ is odd, since
	$\gcd(4,a_2,a_3)=1$.  The integer $p$ is also odd, since otherwise $2p$
	would be a multiple of $4$, contrary to the minimality of the generating
	set.  In particular, $p\geq3$.
	Suppose that $a\leq p$.  Since $a$ and $p$ are odd, we have
	$2p=2a+4(p-a)/2\in\left<4,a\right>$, again contradicting minimality.
	Therefore $a>p$, so $a=p+2q$ for some $q\geq1$.  In particular,
	$a\in\left<2,p\right>$, and Theorem~\ref{thm:herzog}(1) shows that $H$ is
	symmetric.
\end{proof}


\subsection{Strategy for Cohen--Macaulayness}
\label{subsec:cm-strategy}


We next explain the strategy used to prove the Cohen--Macaulayness of the Rees algebra.


For a commutative ring $A$ and an ideal $J$ of $A$, put
$$
\calR_A(J)=A[JT]=\bigoplus_{n\geq0}J^nT^n\subseteq A[T], \qquad \gr_J(A)=\bigoplus_{n\geq0}J^n/J^{n+1}.
$$
where $T$ is an indeterminate over $A$.

Throughout this subsection, let $S=k[x,y,z]$, let $\m=(x,y,z)$, and let $I$ be an $\m$-primary homogeneous ideal of $S$. 
For simplicity, $\calR_S(I)$ will be denoted simply as $\calR(I)$ in what follows.


Our strategy is the standard one of proving that the reduction number is one.  
However, even when $I$ is a monomial ideal, it can be quite difficult to find a minimal reduction of $I$ generated by $3=\dim S$ elements, even in situations where the existence of such a reduction is clear, for example when $k$ is infinite. The ideals considered in this paper are not monomial.
We therefore seek a three-generated ideal whose localization at $\m$ is a reduction of $IS_\m$, and verify directly that the local reduction number is one.

The following proposition shows that this local calculation suffices to prove that $\calR(I)$ is Cohen--Macaulay.  We state it in three variables, which is the setting of this paper, although the same proof works in $n$ variables.

\begin{proposition}
\label{prop:local-reduction-criterion}
Suppose that there is a three-generated ideal $Q\subseteq I$ such that
\[
 I^2=QI+\m I^2.
\]
Then $\calR(I)$ is Cohen--Macaulay.
\end{proposition}

\begin{proof}
Put $A=S_\m$.  The assumed equality
gives
\[
 (IA)^2=(QA)(IA)+(\m A)(IA)^2,
\]
and hence Nakayama's lemma yields $(IA)^2=(QA)(IA)$.  Thus $QA$ is a
reduction of the $\m A$-primary ideal $IA$.  Since $Q$ is generated by three
elements and $\dim A=3$, the ideal $QA$ is a parameter ideal for $A$ and
hence a minimal reduction of $IA$.
By the Valabrega--Valla criterion \cite[Proposition~3.1]{ValabregaValla1978}, the
associated graded ring $\gr_{IA}(A)$ is Cohen--Macaulay.  Since the reduction
number of $IA$ with respect to $QA$ is at most one, it is at most
$\dim A-1=2$.  The Goto--Shimoda criterion
\cite[Theorem~3.1]{GotoShimoda1982} now shows
that $\calR_A(IA)$ is Cohen--Macaulay.

It remains to pass back to $S$.  Let $P\in\Spec\calR(I)$ and put
$\fkp=P\cap S$.  The ring $\calR(I)_P$ is a localization of
$\calR_{S_\fkp}(IS_\fkp)$.  If $\fkp=\m$, $\calR(I)_P$ is a localization of the
Cohen--Macaulay ring $\calR_A(IA)$.  If $\fkp\ne\m$, then
$IS_\fkp=S_\fkp$, since $\sqrt I=\m$.  Hence
\[
 \calR_{S_\fkp}(IS_\fkp)=S_\fkp[T],
\]
which is Cohen--Macaulay.  Thus $\calR(I)$ is Cohen--Macaulay.
\end{proof}


To prove the Cohen--Macaulay part of Theorem~\ref{thm:main}, we give a three-generated ideal $Q\subseteq I$ for each ideal $I$ under consideration.
After writing
\[
 I=Q+(f_1,f_2,\ldots,f_r),
\]
it is enough to check that $f_if_j\in QI+\m I^2$ for every $1\leq i\leq j\leq r$.  These inclusions give $I^2=QI+\m I^2$, and Proposition~\ref{prop:local-reduction-criterion} then shows that $\calR(I)$ is Cohen--Macaulay.


\subsection{Strategy for Normality}
\label{subsec:normality}

We now describe our strategy for proving the normality of the ideals in the families under consideration.


Let $S=k[x_1,\ldots,x_n]$ have a positive weighted grading, let
$\m=(x_1,\ldots,x_n)$, and let $I=(g_1,\ldots,g_\mu)$ be an
$\m$-primary homogeneous ideal with this minimal system of generators;
thus $\mu=\mu_S(I)$. Suppose that there is a homogeneous element
$f$ of degree $\ell$ such that $J=I+(f)$ is normal and $\m f\subseteq I$.
In the applications below, $g_1=f-f'$ is a binomial with monomial terms
$f,f'$, and $J=I+(f)=I+(f')$ is an integrally closed monomial ideal
with at most seven generators. Its normality follows from
\cite{AtakaMatsuoka2026}; the inclusion $\m f\subseteq I$ is checked
for each family. The argument below does not depend on the number of generators of $I$.

Suppose, to the contrary, that $I^s$ is not integrally closed for some $s\geq1$.
Since the integral closure of a homogeneous ideal is again homogeneous, 
we may take a homogeneous element
$\xi\in\ol{I^s}\setminus I^s$
of degree $\delta$. 
Then the normality of $J$ gives
$$
 \xi\in \ol{I^s} \subseteq \ol{J^s} = J^s
 = I^s + fI^{s-1} + \cdots + (f^s).
$$
Now choose the least positive integer $m$ such that
$$
 \xi\in I^s+fI^{s-1}+\cdots+f^m I^{s-m}.
$$
We then have
$$
 \xi=\xi_0+f^m\omega,
$$
for some $\xi_0\in \fka := I^s+fI^{s-1}+\cdots+f^{m-1}I^{s-m+1}$ and $\omega\in I^{s-m}$.
Since $\omega$ is homogeneous of degree $\delta-m\ell$, we may choose homogeneous elements $c_\lambda\in S$, indexed by $\lambda=(\lambda_1,\lambda_2,\ldots,\lambda_{\mu})\in\mathbb N^{\mu}$ with $\lambda_1+\lambda_2+\cdots+\lambda_{\mu}=s-m$.  For such a multi-index, write $g^\lambda=g_1^{\lambda_1}g_2^{\lambda_2}\cdots g_{\mu}^{\lambda_{\mu}}$.  Then
$$
\omega=\sum_{\lambda}c_\lambda g^\lambda,
$$
where
$$
\deg c_\lambda+\lambda_1\deg g_1+\lambda_2\deg g_2+\cdots+\lambda_{\mu}\deg g_{\mu}=\delta-m\ell
$$
whenever $c_\lambda\ne0$.  Since $\m f\subseteq I$, we have
$$
f^m\m I^{s-m}\subseteq f^{m-1}I^{s-m+1}\subseteq\fka.
$$
Thus, after absorbing into $\xi_0$ all terms for which $c_\lambda\in\m$, we may assume that every $c_\lambda$ belongs to $k$.  Consequently, only products of weighted degree $\delta-m\ell$ remain.  For $p,q\in\mathbb N$, put
$$
\begin{aligned}
\Lambda^p_q
=\{\lambda = (\lambda_1,\lambda_2,\ldots,\lambda_{\mu})\in\mathbb N^{\mu}\mid {}&
\lambda_1+\lambda_2+\cdots+\lambda_{\mu}=p,\\
&\lambda_1\deg g_1+\lambda_2\deg g_2+\cdots+\lambda_{\mu}\deg g_{\mu}=q\}.
\end{aligned}
$$
We may now write
$$
\omega=\sum_{\lambda\in\Lambda^{s-m}_{\delta-m\ell}}c_\lambda g^\lambda.
$$

If $fg_i\in I^2$ for some $i$, then, for every $\lambda\in\Lambda^{s-m}_{\delta-m\ell}$ with $\lambda_i>0$, we have $c_\lambda f^m g^\lambda\in\fka$.  By absorbing this term into $\xi_0$, we may therefore replace $c_\lambda$ by zero.
Similarly, if $fg_ig_j\in I^3$ for some $i\ne j$, we may take $c_\lambda=0$ whenever $\lambda_i>0$ and $\lambda_j>0$.  In this way, we restrict the set over which $\lambda$ ranges in each case and denote the resulting subset of $\Lambda^{s-m}_{\delta-m\ell}$ by $\Lambda$.  Thus
$$
\omega = \sum_{\lambda \in \Lambda} c_\lambda g^\lambda.
$$
More generally, suppose that two products $g^\lambda$ and $g^{\lambda'}$
have the same number of factors and the same weighted degree.  If
$g^\lambda-g^{\lambda'}\in\m I^{s-m}$, then
$f^m(g^\lambda-g^{\lambda'})\in\fka$, so either product may be replaced by
the other in $\omega$ without changing $f^m\omega$ modulo $\fka$.
The preceding reductions are not exhaustive.  In some cases, further adjustments of the exponent vectors are made using relations among products of the generators.  All such adjustments, together with the resulting set $\Lambda$, will be specified explicitly in each case.


The remaining step is to separate the coefficients.  We first fix the data
arising from the preceding reduction.  Let $k$ be an infinite field, let
$n,\mu\geq1$ be integers, and let $S=k[x_1,x_2,\ldots,x_n]$.  Let
$I=(g_1,g_2,\ldots,g_\mu)$ be an ideal of $S$, and let $f\in S$.  For
$\lambda=(\lambda_1,\lambda_2,\ldots,\lambda_\mu)\in\mathbb N^\mu$, write
$|\lambda|=\lambda_1+\lambda_2+\cdots+\lambda_\mu$.
Choose integers $s,m$ with $1\leq m\leq s$, a subset
$\Lambda\subseteq\{\lambda\in\mathbb N^\mu\mid |\lambda|=s-m\}$,
an element $\xi_0\in\sum_{j=0}^{m-1}f^jI^{s-j}$, and coefficients
$c_\lambda\in k$ for $\lambda\in\Lambda$.  Put
$$
 \omega=\sum_{\lambda\in\Lambda}c_\lambda g^\lambda,
 \qquad \xi=\xi_0+f^m\omega,
$$
and suppose that $\xi\in\ol{I^s}$.

The essential step is to construct one-variable maps that separate the
remaining products.  We record the resulting criterion.

\begin{lemma}\label{lem:coefficient-separation}
Let $a$ and $\nu$ be positive integers.
For $\varepsilon=(\varepsilon_1,\varepsilon_2,\ldots,\varepsilon_a)
\in(k^\times)^a$ and $\alpha=(\alpha_1,\alpha_2,\ldots,\alpha_a)
\in\mathbb Z^a$, write
$\varepsilon^\alpha=\varepsilon_1^{\alpha_1}\varepsilon_2^{\alpha_2}
\cdots\varepsilon_a^{\alpha_a}$.
For each $\varepsilon\in(k^\times)^a$, suppose that there is a
$k$-algebra homomorphism $\psi_\varepsilon:S\to k[[t]]$ such that
\[
 \ord_t\psi_\varepsilon(f)=\nu-1,\qquad
 \ord_t\psi_\varepsilon(g_i)\geq\nu\quad(1\leq i\leq\mu).
\]
Here, $\ord_t$ denotes the discrete valuation of $k[[t]]$.
Suppose that the set
\[
 G_\varepsilon=\{i\in\{1,\ldots,\mu\}\mid
 \ord_t\psi_\varepsilon(g_i)=\nu\}
\]
is independent of $\varepsilon$, and denote this common set by $G$.
Suppose further that, for each $i\in G$, there are
$u_i\in k^\times$ and $\alpha_i\in\mathbb Z^a$ such that, for every
$\varepsilon\in(k^\times)^a$,
\[
 \psi_\varepsilon(g_i)
 =u_i\varepsilon^{\alpha_i}t^\nu+\text{(higher-order terms)}.
\]
For $\lambda\in\mathbb N^\mu$, write
\[
 \operatorname{supp}\lambda
 =\{i\in\{1,\ldots,\mu\}\mid\lambda_i>0\}.
\]
Put $\Lambda_G=\{\lambda\in\Lambda\mid\operatorname{supp}\lambda\subseteq G\}$.
If $\lambda\mapsto\sum_{i\in G}\lambda_i\alpha_i$ is injective on
$\Lambda_G$, then $c_\lambda=0$ for every $\lambda\in\Lambda_G$.
\end{lemma}

\begin{proof}
Applying $\psi_\varepsilon$ to an equation of integral dependence for
$\xi$ shows that $\psi_\varepsilon(\xi)$ is integral over
$\psi_\varepsilon(I^s)k[[t]]$. Since every ideal of the discrete valuation
ring $k[[t]]$ is integrally closed, we have
\[
 \psi_\varepsilon(\xi)\in\psi_\varepsilon(I^s)k[[t]]
 \subseteq(t^{s\nu}).
\]
On the other hand, since
 $
 \ord_t\psi_\varepsilon(\xi_0)\geq s\nu-m+1.
 $
As $m\geq1$, the equality $\xi=\xi_0+f^m\omega$ therefore gives
\[
 \ord_t\psi_\varepsilon(\omega)\geq(s-m)\nu+1.
\]
The coefficient of $t^{(s-m)\nu}$ in $\psi_\varepsilon(\omega)$ is thus zero:
\[
 \sum_{\lambda\in\Lambda_G}
 c_\lambda\left(\prod_{i\in G}u_i^{\lambda_i}\right)
 \varepsilon^{\sum_{i\in G}\lambda_i\alpha_i}=0.
\]
This holds for every $\varepsilon\in(k^\times)^a$, so it is an identity
of Laurent polynomials over the infinite field $k$. Injectivity makes
the parameter monomials distinct; since each $u_i$ is nonzero, the
asserted coefficients vanish.
\end{proof}

In our applications, $\Lambda$ also has fixed weighted degree
$\sum_i\lambda_i\deg g_i=\delta-m\ell$. The parameter exponents,
together with this equality and $|\lambda|=s-m$, determine $\lambda$;
we verify this injectivity separately in each case.
If one map has $\Lambda_G=\Lambda$, the lemma gives $\omega=0$,
contrary to the minimality of $m$. If several maps are needed, we apply
the lemma successively to the remaining coefficients and check that
the subsets treated by these maps cover $\Lambda$.
